\documentclass[12pt]{letter}
\usepackage[paperwidth=8.5in,paperheight=11in,margin=1.1in]{geometry}
\usepackage{mathptmx}
\usepackage{hyperref}
\pagestyle{empty}
\signature{Guanyu Wang, M.D., Ph.D.\\
Department of General Surgery\\
Sir Run Run Shaw Hospital, School of Medicine\\
Zhejiang University, Hangzhou 310016, China\\
Email: wangguanyu@zju.edu.cn}

\begin{document}

\begin{letter}{The Editors\\
\textit{Frontiers in Oncology}}

\opening{Dear Editors,}

We are pleased to submit our manuscript entitled ``Identification of Long Non-Coding RNA Signatures for Diagnosis and Prognosis of Colorectal Cancer via Machine Learning'' for consideration for publication in \textit{Frontiers in Oncology}.

In this study, we developed an integrated computational pipeline to identify lncRNA signatures with both diagnostic and prognostic value in colorectal cancer (CRC). Using harmonized TCGA--GTEx expression data (639 samples) and TCGA-COAD clinical follow-up data (428 patients), we identified an 8-lncRNA panel whose expression profile, together with pathological stage, enabled robust risk stratification and accurate discrimination of tumor from normal tissue. Key findings include: (i) a Cox model incorporating eight lncRNAs and stage stratified patients into distinct risk groups (log-rank $p < 0.0001$), with time-dependent AUC converging to 0.725 at 8 years; (ii) SVM, random forest, and neural network classifiers achieved testing accuracies of 91.6\%, 92.2\%, and 90.6\%, substantially outperforming baseline methods; and (iii) several lncRNAs in our panel (EIF3J-AS1, MIR210HG, LINC00261) have been independently validated in prior studies, while others represent potentially novel CRC biomarkers.

The manuscript has not been previously published and is not under consideration elsewhere. All analyses use publicly available datasets under appropriate guidelines. Complete R scripts are available at \url{https://github.com/woodhaha/CRC_data_mining} to ensure reproducibility.

We believe this work will interest the readership of \textit{Frontiers in Oncology} as a systematic, reproducible framework for transcriptomic biomarker discovery. We appreciate your consideration and look forward to hearing from you.

\closing{Sincerely,}

\end{letter}
\end{document}
